\section{Proof of Proposition~\texorpdfstring{\ref{prop:aging-expansion}}{5.5}}\label{app:aging}

We use the setting of Section~\ref{sec:aging} with two states and a fair first
step. So $0<c<2$, $X_1$ is a fair sign, and the switch indicators
$\chi_k=\mathbf 1\{X_k\ne X_{k-1}\}$, $k\ge2$, are independent
Bernoulli$(c/k)$ variables, independent of $X_1$. Put
\[
\rho(a,b)=\prod_{k=a+1}^{b}\Bigl(1-\frac ck\Bigr)
=\frac{\Gamma(b+1-c)\,\Gamma(a+1)}{\Gamma(b+1)\,\Gamma(a+1-c)},\qquad
\lambda_{t,N}=\frac ct\,\rho(t,N)
\]
for integers $1\le a\le b$ and $2\le t\le N$. So $\lambda_{t,N}$ is the
probability of a switch at time $t$ and none in $(t,N]$, and
$E_N=\frac12\rho(1,N)$. Let $T^+_s$ be the time of the $s$-th step equal to
$+1$. Let $B$ be a Beta$(c,c)$ variable, $d_c=4^{1-c}/\mathrm B(c,c)$ its
density at $\frac12$, and $\Theta(c)=\log\bigl(2\Gamma(2-c)\,d_c\bigr)$. For
even $N$ we write $h=N/2$ and $C_N=\Prob(A_N=h)$. Changing the sign of every
step preserves the law of the walk and sends $A_N$ to $N-A_N$. We call this
flip symmetry.

\paragraph{Idea of the proof.}
A walk that ends in the center with last step $-1$ takes its $h$-th step
$+1$ at some time $r$, switches at time $r+1$, and never switches again. So
$C_N$ is an average of the explicit numbers $\lambda_{r+1,N}$ over the law of
$T^+_h$ (Lemma~\ref{lem:D-average}). For $c\le1$ these numbers are monotone in
$r$, which already traps $c_N$ between two explicit roots and gives
$1-c_N=O(1/\log N)$ (Proposition~\ref{prop:D-bracket}). For the constant we
need the law of $T^+_h$. We have $T^+_h\le r$ exactly when $A_r\ge h$, and
$A_r/r$ is close to a Beta$(c,c)$ variable with an error that is a power of
$r$ (Lemma~\ref{lem:D-coupling}). This is proved by coupling the walk with a
sign process in continuous time that switches at the points of a Poisson
process of intensity $c\,dt/t$, whose occupation fraction is exactly
Beta$(c,c)$ at every time (Proposition~\ref{prop:D-beta}). Summation by parts
then gives $NC_N=d_c+O(N^{-0.4})$ (Proposition~\ref{prop:D-llt}).

\subsection{A gamma ratio}

Gautschi's inequality~\cite[5.6.4]{dlmf} says that
$x^{1-s}<\Gamma(x+1)/\Gamma(x+s)<(x+1)^{1-s}$ for $x>0$ and $0<s<1$. We
claim that for $0<c<2$ and real $y\ge2$,
\begin{equation}\label{eq:D-gamma}
\frac{\Gamma(y+1-c)}{\Gamma(y+1)}=y^{-c}(1+\vartheta)\qquad\text{with}\qquad
-\frac1y\le\vartheta\le\frac2y .
\end{equation}
If $c=1$ the left side is $1/y$ and $\vartheta=0$. If $c<1$ we take $x=y$ and
$s=1-c$. Then the left side lies between $(y+1)^{-c}$ and $y^{-c}$, and
$(1+1/y)^{-c}\ge1-c/y\ge1-1/y$. If $c>1$ we take $x=y$ and $s=2-c$ and use
$\Gamma(y+1-c)=\Gamma(y+2-c)/(y+1-c)$. Then the left side lies between
$(y+1)^{1-c}/(y+1-c)$ and $y^{1-c}/(y+1-c)$. Here
$y/(y+1-c)=1+(c-1)/(y+1-c)\le1+1/(y-1)\le1+2/y$, and
$(1+1/y)^{1-c}\,y/(y+1-c)\ge(1+1/y)^{-1}\ge1-1/y$.

\subsection{Averaging at the center and a bracket}

\begin{lemma}[averaging at the center]\label{lem:D-average}
Let $N\ge2$ be even and $T=T^+_h$. Then
\begin{equation}\label{eq:D-average}
C_N=2\sum_{r=h}^{N-1}\Prob(T=r)\,\lambda_{r+1,N}\qquad\text{and}\qquad
\sum_{r=h}^{N-1}\Prob(T=r)=\frac12 .
\end{equation}
If $0<c\le1$, then $c/N\le C_N\le\lambda_{h+1,N}$.
\end{lemma}

\begin{proof}
On the event $\{A_N=h,\,X_N=-1\}$ the last step equal to $+1$ is the $h$-th
one, at a time $T=r\le N-1$. It is followed by a switch at $r+1$ and no switch
in $(r+1,N]$. Conversely, these conditions give $A_N=h$ and $X_N=-1$. The
event $\{T=r\}$ depends only on $X_1,\chi_2,\dots,\chi_r$. So
$\Prob(A_N=h,\,X_N=-1)=\sum_{r=h}^{N-1}\Prob(T=r)\,\lambda_{r+1,N}$, and flip
symmetry gives the same value for $X_N=+1$. Next, $T\le N-1$ exactly when
$A_{N-1}\ge h$. This has probability $\frac12$, because $N-1$ is odd and
$A_{N-1}$ is symmetric about $(N-1)/2$. Finally
$\lambda_{t+1,N}/\lambda_{t,N}=t/(t+1-c)$, which is at most $1$ for $c\le1$.
So $\lambda_{N,N}=c/N\le\lambda_{t,N}\le\lambda_{h+1,N}$ for
$h+1\le t\le N$, and by~\eqref{eq:D-average} $C_N$ is a weighted average of
these values.
\end{proof}

For an integer $M\ge2$ put $\eta_M(c)=M\rho(1,M)/(2c)$. Since $\rho(1,M)$
decreases in $c$ and equals $1/M$ at $c=1$, the function $\eta_M$ decreases
strictly on $(0,1]$, from $\infty$ at $0$ to $\eta_M(1)=\frac12$. Let
$c^{(M)}\in(0,1)$ be the root of $\eta_M=1$.

\begin{proposition}[bracket]\label{prop:D-bracket}
The numbers $c^{(M)}$ increase with $M$, and
$1-c^{(M)}=\frac{\log2}{\log M+1+\gamma}+O\bigl((\log M)^{-3}\bigr)$. For
every even $N\ge2$ we have $c^{(h+1)}\le c_N\le c^{(N)}$. So
$(1-c_N)\log N\to\log2$, and $c_N\ge c^{(3)}=(9-\sqrt{57})/2>0.725$ for
$N\ge4$.
\end{proposition}

\begin{proof}
For $c<1$ we have $\eta_{M+1}/\eta_M=(M+1-c)/M>1$. So
$\eta_{M+1}(c^{(M)})>1$ and $c^{(M+1)}>c^{(M)}$. By
Proposition~\ref{prop:aging-fair}(b), $C_N<E_N$ for $c<c_N$ and $C_N>E_N$ for
$c>c_N$, and $c_N<1$. So it is enough to compare $C_N$ with $E_N$ for
$0<c<1$. If $c^{(N)}<c<1$, then $\eta_N(c)<1$, that is
$E_N=\frac12\rho(1,N)<c/N\le C_N$ by Lemma~\ref{lem:D-average}. So
$c_N\le c^{(N)}$. If $c<c^{(h+1)}$, then $\eta_{h+1}(c)>1$. With
$\rho(1,N)=\rho(1,h+1)\rho(h+1,N)$ this gives
$E_N>\frac{c}{h+1}\rho(h+1,N)=\lambda_{h+1,N}\ge C_N$. So
$c_N\ge c^{(h+1)}$.

The number $c^{(3)}$ solves $3\rho(1,3)=(2-c)(3-c)/2=2c$, that is
$c^2-9c+6=0$. By~\eqref{eq:D-gamma} with $y=M$, the equation $\eta_M(c)=1$
reads $(1-c)\log M=\Theta_0(c)+O(1/M)$ with
$\Theta_0(c)=\log\bigl(2c\,\Gamma(2-c)\bigr)$, uniformly for $0<c\le1$. For
$M\ge3$ we have $c^{(M)}\in[c^{(3)},1)$, where $\Theta_0$ is bounded. So
$e=1-c^{(M)}$ is $O(1/\log M)$, and we may expand at $c=1$. We have
$\Theta_0(1)=\log2$ and $\Theta_0'(1)=1-\psi(1)=1+\gamma$, where
$\psi=\Gamma'/\Gamma$. Taylor's theorem gives
$e(\log M+1+\gamma)=\log2+O(e^2)+O(1/M)$, which is the expansion. The limit
of $(1-c_N)\log N$ follows since $\log(h+1)\sim\log N$.
\end{proof}

\subsection{The continuum occupation}

\begin{proposition}[continuum occupation]\label{prop:D-beta}
Let $\Pi$ be a Poisson process on $(0,\infty)$ with intensity $c\,dt/t$, and
let $Y(1)$ be a fair sign independent of $\Pi$. Put
$Y(t)=Y(1)(-1)^{\Pi((1,t])}$ for $t\ge1$ and $Y(t)=Y(1)(-1)^{\Pi((t,1])}$ for
$0<t<1$. Then for every $r>0$ the fraction
$Z_r=\frac1r\int_0^r\mathbf 1\{Y(t)=+1\}\,dt$ has exactly the Beta$(c,c)$ law.
\end{proposition}

This agrees with~\cite[Remark~6.2]{evw2020}, which obtains the Beta$(c,c)$
one-dimensional marginals of the zigzag process
of~\cite[Definition~6.1]{evw2020} from a scaling limit and mentions a direct
proof by E.~Crane that is not written there. We give a direct proof.

\begin{proof}
The points $\log t$, $t\in\Pi$, form a Poisson process $\tilde\Pi$ of rate $c$
on $\mathbb R$, and $W(u)=Y(e^u)$ satisfies
$W(v)=W(u)(-1)^{\tilde\Pi((u,v])}$ for $u<v$. For each real $a$, $W(a)$ is
$W(0)$ times a function of $\tilde\Pi$, so it is a fair sign independent of
$\tilde\Pi$. As $\tilde\Pi-a$ has the law of $\tilde\Pi$, the pair
$(W(a),\tilde\Pi-a)$ has the law of $(W(0),\tilde\Pi)$, and so $W(a+\cdot)$
has the law of $W$. Hence $Z_r=\int_0^1\mathbf 1\{Y(rt)=+1\}\,dt$ has the law
of $Z_1=\int_0^\infty e^{-s}\mathbf 1\{V_s=+1\}\,ds$, where $V_s=W(-s)$ starts
at a fair sign and switches at rate $c$. Let $\mu_\pm$ be the law of $Z_1$
given $V_0=\pm1$, let $\tau$ be the first switch of $V$, and let
$U=e^{-\tau}$. Then $\Prob(U\le u)=u^c$ for $0\le u\le1$, so $U$ is
Beta$(c,1)$. After $\tau$ the process starts afresh from $-V_0$,
independently of $\tau$. Splitting the integral at $\tau$ shows that
$(\mu_+,\mu_-)$ is a fixed point of the map $\Psi$ that sends a pair
$(\nu_+,\nu_-)$ of laws on $[0,1]$ to the laws of $1-U+UZ_-$ and $UZ_+$,
where $Z_\pm\sim\nu_\pm$ are independent of $U$. Coupling $Z_\pm$ optimally
with $Z'_\pm\sim\nu'_\pm$, independently of $U$, gives
$\E\lvert UZ_\pm-UZ'_\pm\rvert=\E U\,\E\lvert Z_\pm-Z'_\pm\rvert$. So $\Psi$
contracts the larger of the two Wasserstein distances by the factor
$\E U=c/(c+1)$, and it has only one fixed point.

Let $G_a,G_b,G_d$ be independent Gamma variables with shapes $a,b,d$ and unit
scale. A change of variables shows that $G_a/(G_a+G_b)$ is Beta$(a,b)$ and
independent of $G_a+G_b$. Hence it is independent of the pair
$(G_a+G_b,G_d)$ and of $(G_a+G_b)/(G_a+G_b+G_d)$, which is Beta$(a+b,d)$.
Their product $G_a/(G_a+G_b+G_d)$ is Beta$(a,b+d)$. With $(a,b,d)=(c,1,c)$
this says that $UZ$ is Beta$(c,c+1)$ whenever $Z$ is Beta$(c+1,c)$ and
independent of $U$. Applying it to $Z=Z_+$ and to $Z=1-Z_-$, and using
$1-U+UZ_-=1-U(1-Z_-)$, we see that $\Psi$ maps the pair Beta$(c+1,c)$,
Beta$(c,c+1)$ to itself. So $Z_1$ is the equal mixture of these two laws. As
$\mathrm B(c+1,c)=\mathrm B(c,c+1)=\frac12\mathrm B(c,c)$, the mixture has
density $x^{c-1}(1-x)^{c-1}(x+1-x)/\mathrm B(c,c)$, which is the Beta$(c,c)$
density.
\end{proof}

\begin{lemma}[coupling]\label{lem:D-coupling}
Let $J\subset(0,2)$ be compact, $a_J=\min(\min J,1)$ and
$\kappa_J=a_J/(1+a_J)$. There is $C_J$ such that for $c\in J$ and every
integer $r\ge2$,
\[
\sup_{y\in\mathbb R}\bigl\lvert\Prob(A_r\ge y)-\Prob(rB\ge y)\bigr\rvert
\le C_J\,(\log r)\,r^{-\kappa_J}.
\]
\end{lemma}

\begin{proof}
Take $\Pi$ and $Y$ as in Proposition~\ref{prop:D-beta} and an integer
$1\le K\le r$. The variables
$\chi'_k=\mathbf 1\{\Pi((k-1,k])\text{ is odd}\}$, $k>K$, are independent of
each other and of $Y(K)$, and
$\Prob(\chi'_k=1)=p'_k=\frac12\bigl(1-(1-1/k)^{2c}\bigr)$. By Taylor's theorem
$\lvert p'_k-c/k\rvert=\frac12c\lvert2c-1\rvert(1-\vartheta)^{2c-2}k^{-2}$ for
some $0<\vartheta<1/k$. This is at most $3k^{-2}$. Indeed, if $c\ge1$ then
$\frac12c\lvert2c-1\rvert<3$ and the power of $1-\vartheta$ is at most $1$.
If $c<1$ then $\frac12c\lvert2c-1\rvert\le\frac12$ and the power is at most
$(1-\frac12)^{-2}=4$, as $k\ge2$. For $K<k\le r$ we couple
$\chi_k\sim\text{Bernoulli}(c/k)$ with $\chi'_k$ so that
$\Prob(\chi_k\ne\chi'_k)=\lvert p'_k-c/k\rvert$, using independent extra
randomness. We put $X_K=Y(K)$ and $X_k=X_{k-1}(-1)^{\chi_k}$ for $K<k\le r$,
and we draw $X_1,\dots,X_{K-1}$ from the law of the walk given $X_K$, again
with independent randomness. In the walk $X_K$ is a fair sign and
$(\chi_k)_{k>K}$ is independent of $(X_1,\dots,X_K)$. So
$(X_1,\dots,X_r)$ has the law of the walk.

Let $\mathcal A$ be the event that $\chi_k=\chi'_k$ for $K<k\le r$. On
$\mathcal A$ we have $X_k=Y(k)$ for $K\le k\le r$, and
$\Prob(\mathcal A^c)\le\sum_{k>K}3k^{-2}\le3/K$. If also $\Pi((k-1,k])=0$,
then $Y=Y(k)$ on $(k-1,k]$, so step $k$ adds the same amount to $A_r$ as the
interval $(k-1,k]$ adds to $rZ_r$. The other steps $k\le r$ number at most
$K+\Pi((K,r])$, and each changes the difference by at most $1$. Hence
$\lvert A_r-rZ_r\rvert\le K+\Pi((K,r])$ on $\mathcal A$. The variable
$\Pi((K,r])$ is Poisson with mean $c\log(r/K)\le2\log r$, so by Markov's
inequality it exceeds $K$ with probability at most $2\log r/K$. With
$b_K=(3+2\log r)/K$ we get
\[
\Prob(rZ_r\ge y+2K)-b_K\le\Prob(A_r\ge y)\le\Prob(rZ_r\ge y-2K)+b_K ,
\]
and $rZ_r$ has the law of $rB$. So the left side of the lemma is at most
$b_K+\sup_x\Prob(x\le B\le x+2K/r)$.

For $c\le1$ the density of $B$ is at most
$2^{1-c}\bigl(x^{c-1}+(1-x)^{c-1}\bigr)/\mathrm B(c,c)$, and $x^{c-1}$ has
integral at most $\delta^c/c$ over an interval of length $\delta$. For
$c\ge1$ the density is at most $d_c$. The mass of an interval is also at
most $1$. So there is $C'_J$ such that an interval of length $\delta$ has
mass at most $C'_J\delta^{a_J}$ for all $\delta>0$ and $c\in J$. Take
$K=\lceil r^{\kappa_J}\rceil$. Then $b_K\le(3+2\log r)r^{-\kappa_J}$ and
$2K/r\le4r^{\kappa_J-1}$, and $(r^{\kappa_J-1})^{a_J}=r^{-\kappa_J}$ because
$(1-\kappa_J)a_J=\kappa_J$.
\end{proof}

\subsection{A local limit at the center}

\begin{proposition}[local limit at the center]\label{prop:D-llt}
For each compact $J\subset(0,2)$ there is $C_J$ such that
$\bigl\lvert NC_N-d_c\bigr\rvert\le C_J\,(\log N)\,N^{-\kappa_J}$ for $c\in J$
and even $N\ge4$.
\end{proposition}

\begin{proof}
Let $T=T^+_h$, $F(r)=\Prob(T\le r)=\Prob(A_r\ge h)$, $\Phi(r)=\Prob(rB\ge h)$
and $\Lambda_r=N\lambda_{r+1,N}$. Then $F(h-1)=\Phi(h-1)=\Phi(h)=0$. For
$h\le r\le N-1$ we have $2\le r$ and $N/2\le r<N$, so
Lemma~\ref{lem:D-coupling} gives
$\lvert F(r)-\Phi(r)\rvert\le\delta_N:=C_J(\log N)N^{-\kappa_J}$ after
enlarging $C_J$. By Lemma~\ref{lem:D-average},
$NC_N=2\sum_{r=h}^{N-1}\bigl(F(r)-F(r-1)\bigr)\Lambda_r$. The ratio
$\lambda_{t+1,N}/\lambda_{t,N}=t/(t+1-c)$ shows that $\Lambda_r$ is monotone
in $r$. Also $\Lambda_{N-1}=c$ and $0<\Lambda_h<2c$. For $c\le1$ this holds
because $\rho\le1$ and $N<2(h+1)$, and for $c>1$ because $\Lambda_r$
increases. So summation by parts shows that replacing $F$ by $\Phi$ changes
$NC_N$ by at most
$2\delta_N\bigl(\Lambda_{N-1}+\lvert\Lambda_{N-1}-\Lambda_h\rvert\bigr)\le4c\,\delta_N$.
For $h<r\le N-1$ we have $\Phi(r)-\Phi(r-1)=\Prob(h/r\le B<h/(r-1))$, and on
this event $\lceil h/B\rceil=r$. So the new main term is
$2\,\E\bigl[\Lambda_{\lceil h/B\rceil};\,B\ge h/(N-1)\bigr]$.

By~\eqref{eq:D-gamma} with $y=N$ and $y=r+1\ge3$,
$\Lambda_r=c\bigl((r+1)/N\bigr)^{c-1}(1+O(1/N))$ for $h\le r\le N-1$. On the
event above with $r=\lceil h/B\rceil$, the number $(r+1)/N$ lies in
$[\frac12,2]$ and within $2/N$ of $1/(2B)$, which also lies in
$[\frac12,2]$. The function $x^{c-1}$ is $4$-Lipschitz on $[\frac12,2]$. So
$\Lambda_{\lceil h/B\rceil}=c(2B)^{1-c}+O(1/N)$. Also
$\Prob\bigl(\frac12\le B<h/(N-1)\bigr)=O_J(1/N)$, since the density of $B$ is
bounded near $\frac12$ uniformly on $J$. So $NC_N$ is within $4c\,\delta_N$ of
\[
2c\,\E\bigl[(2B)^{1-c};\,B\ge\tfrac12\bigr]+O_J(1/N)
=\frac{2c\,2^{1-c}}{\mathrm B(c,c)}\int_{1/2}^1(1-b)^{c-1}\,db+O_J(1/N)=d_c+O_J(1/N).\qedhere
\]
\end{proof}

\subsection{The expansion of the crossing}

\begin{proof}[Proof of Proposition~\ref{prop:aging-expansion}]
Let $J=[c^{(3)},1]$. Then $\kappa_J=c^{(3)}/(1+c^{(3)})>0.42$, and $d_c$ is
bounded below on $J$ by a positive constant. By~\eqref{eq:D-gamma} with
$y=N$,
$\log E_N=\log\bigl(\tfrac12\rho(1,N)\bigr)=-\log2-\log\Gamma(2-c)-c\log N+O(1/N)$.
By Proposition~\ref{prop:D-llt},
$\log C_N=\log d_c-\log N+O(N^{-0.4})$ for large $N$. Both hold uniformly for
$c\in J$, so
\begin{equation}\label{eq:D-logratio}
\log C_N-\log E_N=(c-1)\log N+\Theta(c)+O(N^{-0.4})\qquad(c\in J).
\end{equation}
Here $\Theta(1)=\log2$ because $d_1=1$. Next,
$\Theta'(c)=-\psi(2-c)-\log4-2\psi(c)+2\psi(2c)$, and with $\psi(1)=-\gamma$
and $\psi(2)=1-\gamma$ this gives $\Theta'(1)=\gamma+2-2\log2$. Put
$M=\log N+\Theta'(1)$. By Proposition~\ref{prop:D-bracket}, $c_N\in J$ for
$N\ge4$ and $e=1-c_N$ is $O(1/\log N)$. The left side
of~\eqref{eq:D-logratio} vanishes at $c=c_N$. As $\Theta$ is smooth on $J$,
Taylor's theorem gives $eM=\log2+O(e^2)+O(N^{-0.4})=\log2+O(M^{-2})$. So
$e=\log2/M+O(M^{-3})$.
\end{proof}

With one more term in Taylor's theorem and $e^2=(\log2)^2/M^2+O(M^{-4})$, the
same argument adds the term $-\Theta''(1)(\log2)^2/(2M^3)$ to the expansion
of $c_N$, with error $O(M^{-4})$. Here $\Theta''(1)=\pi^2/2-4$, since
$\Theta''(c)=\psi'(2-c)-2\psi'(c)+4\psi'(2c)$, $\psi'(1)=\zeta(2)$ and
$\psi'(2)=\zeta(2)-1$. At $N=200$ the bracket of
Proposition~\ref{prop:D-bracket} is $[0.88749,0.89878]$,
Proposition~\ref{prop:aging-expansion} gives $0.89319$, and with the $M^{-3}$
term we get $0.89236$. The high-precision value is $c_{200}=0.89253$. The
constants $C_J$ and the constants in the $O$-terms of this appendix are not
computed.
